% Options for packages loaded elsewhere
\PassOptionsToPackage{unicode}{hyperref}
\PassOptionsToPackage{hyphens}{url}
\documentclass[
]{report}
\usepackage{xcolor}
\usepackage{amsmath,amssymb}
\setcounter{secnumdepth}{5}
\usepackage{iftex}
\ifPDFTeX
  \usepackage[T1]{fontenc}
  \usepackage[utf8]{inputenc}
  \usepackage{textcomp} % provide euro and other symbols
\else % if luatex or xetex
  \usepackage{unicode-math} % this also loads fontspec
  \defaultfontfeatures{Scale=MatchLowercase}
  \defaultfontfeatures[\rmfamily]{Ligatures=TeX,Scale=1}
\fi
\usepackage[]{fontawesome}
\ifPDFTeX\else
  % xetex/luatex font selection
\fi
% Use upquote if available, for straight quotes in verbatim environments
\IfFileExists{upquote.sty}{\usepackage{upquote}}{}
\IfFileExists{microtype.sty}{% use microtype if available
  \usepackage[]{microtype}
  \UseMicrotypeSet[protrusion]{basicmath} % disable protrusion for tt fonts
}{}
\makeatletter
\@ifundefined{KOMAClassName}{% if non-KOMA class
  \IfFileExists{parskip.sty}{%
    \usepackage{parskip}
  }{% else
    \setlength{\parindent}{0pt}
    \setlength{\parskip}{6pt plus 2pt minus 1pt}}
}{% if KOMA class
  \KOMAoptions{parskip=half}}
\makeatother
\setlength{\emergencystretch}{3em} % prevent overfull lines
\providecommand{\tightlist}{%
  \setlength{\itemsep}{0pt}\setlength{\parskip}{0pt}}
\usepackage[explicit]{titlesec}
\usepackage[svgnames]{xcolor}
\usepackage{enumitem}

% section headers

\titleformat{\chapter}
  [hang]
  {\normalfont\sffamily\LARGE\color{white}}
  {}
  {0.5ex}
  {\colorbox{RoyalBlue}{\parbox{\textwidth}{Article \thechapter \hspace{1cm} #1}}}  

% \newcommand{\chapterbreak}{\addpenalty{-500}\vspace*{0pt}}

\titleformat{\section}
  {\normalfont\sffamily}
  {}
  {0.5ex}
  {#1}

% enumerations
\renewcommand{\labelenumi}{\arabic{enumi}}
\renewcommand{\labelenumii}{\arabic{enumi}.\arabic{enumii}}
\renewcommand{\labelenumiii}{\arabic{enumi}.\arabic{enumii}.\arabic{enumiii}}
\renewcommand{\labelenumiv}{\arabic{enumi}.\arabic{enumii}.\arabic{enumiii}.\arabic{enumiv}}

\setlist{noitemstep}

% special environment
\newenvironment{spec}[0]{\footnotesize}{\normalsize}
\usepackage{bookmark}
\IfFileExists{xurl.sty}{\usepackage{xurl}}{} % add URL line breaks if available
\urlstyle{same}
\hypersetup{
  hidelinks,
  pdfcreator={LaTeX via pandoc}}

\author{}
\date{}

\begin{document}

\chapter{Purpose}\label{purpose}

\begin{enumerate}
\tightlist
\item
  The purpose of the Competition Rules is to provide standardized rules
  for all levels of Championships promoted and/or recognized by the
  World Taekwondo Federation (hereafter WT), WT Continental Unions
  (hereafter CUs), or/and WT Member National Associations (hereafter
  MNAs); the Competition Rules is intended to ensure that all matters
  related to competitions are conducted in a fair and orderly manner.
\end{enumerate}

\section{(Interpretation)}\label{interpretation}

\begin{spec}

The objective of Article 1 is to ensure the standardization of all
Taekwondo competition worldwide. Any competition not following the
fundamental principles of these rules cannot be recognized as Taekwondo
competition.

\end{spec}

\chapter{Application}\label{application}

\begin{enumerate}
\item
  The Competition Rules shall apply to all the competitions to be
  promoted and/or recognized by the WT, each CU and MNA. However, any
  MNA wishing to modify some or any part of the Competition Rules must
  first gain the prior approval of the WT. In the case that a CU and/or
  a MNA violates the Competition Rules without prior approval of the WT,
  the WT may exercise its discretion to disapprove or revoke its
  approval of the concerned international tournament. In addition, the
  WT may take further disciplinary actions to the pertinent CU or MNA.
\item
  All competitions promoted or recognized by the WT and /or each CU
  and/or MNA shall observe the WT Statutes, the Bylaws of Dispute
  Resolution and Disciplinary Action, and all other pertinent rules and
  regulations.
\item
  All competitions promoted or recognized by the WT and /or each CU
  and/or MNA shall abide by the WT Medical Code and the WT Anti-Doping
  Rules.
\end{enumerate}

\section{(Explanation\#1)}\label{explanation1}

First gain the approval: Any organization desiring to make a change in
any portion of the existing rules must submit to the WT the contents of
the desired amendment along with the reasons for the desired changes.
Approval for any changes in these rules must be received from the WT at
least one month prior to the scheduled competition. WT can apply
Competition Rules with modifications in its promoted Championships with
the decision of the Technical Delegate after approval of the President.

\chapter{Competition Area}\label{competition-area}

\begin{enumerate}
\item
  The Contest Area shall have a flat surface without any obstructing
  projections, and be covered with an elastic and not slippery mat. The
  Contest Area may also be installed on a platform 0.6-1m high from the
  base, if necessary. The outer part of the Outer Line shall be inclined
  at a gradient of less than 30 degrees, for the safety of the
  contestants. One of the following shapes can be used for Contest Area.

  \begin{enumerate}
  \item
    Square-shape

    The Competition Area is comprised of a Contest Area and a Safety
    Area. The square-shape Contest Area shall be 8m x 8m including 60cm
    boundary line. Surrounding the contest area, approximately
    equidistant on all sides, shall be the Safety Area. The size of
    Competition Area (which envelopes the Contest Area and the Safety
    Area) shall be no smaller than 10m x 10m and no larger than 12m x
    12m. If the Competition Area is on a platform, the Safety Area can
    be increased as needed to ensure the safety of contestants. The
    Contest Area and the Safety Area may be different colors, as
    specified in the relevant competition's Operational Manual.
  \item
    Octagonal-shape

    The Competition Area is comprised of a Contest Area and a Safety
    Area. The Competition Area shall be square shaped and the size shall
    be no smaller than 10mx10m and no larger than 12mx12m. At the center
    of the Competition Areas shall be the octagonal-shape Contest Area.
    The Contest Areas shall measure approximately 8m in diameter, and
    each side of the octagon shall have a length of approximately 3.3m.
    Between the outer line of the Competition area and the boundary line
    of the Contest area is the Safety Area. The Contest Area and the
    Safety Area may be different colors, as specified in the relevant
    competition's Operational Manual.
  \end{enumerate}
\item
  Indication of positions

  \begin{enumerate}
  \item
    The outer line of the Contest Area shall be called the Boundary
    Line(s) and the outer line of the Competition Area shall be called
    the Outer Line(s).
  \item
    The front outer line adjacent to the Recorder's Desk shall be called
    Outer Line \#1, and Clockwise from Outer Line \#1, the other lines
    shall be called Outer Lines \#2, \#3, and \#4. The Boundary Line
    adjacent to the Outer Line \#1 shall be called Boundary Line \#1 and
    clockwise from Boundary Line \#1, the other lines shall be called
    Boundary Lines \#2, \#3, and \#4. In case of Octagonal Shape Contest
    Area, The Boundary Line adjacent to the Outer Line \#1 shall be
    called Boundary Line \#1 and clockwise from Boundary Line \#1, the
    other lines shall be called Boundary Lines \#2, \#3, \#4, \#5, \#6,
    \#7, and \#8.
  \item
    Positions of Referee and Contestants at the beginning and end of the
    Match: The position of the Contestants shall be at the two opposing
    points, 1m from the center point of the Contest Area parallel to
    Outer Line \#1. The Referee shall be positioned at a point 1.5m from
    the center of the Contest Area toward the Outer Line \#3.
  \item
    Positions of Judges: The position of the 1st Judge shall be located
    at a point in minimum 2 m from the corner of the Boundary Line \#2.
    The position of the 2nd Judge shall be located at a point in minimum
    2 m outward from the center of the Boundary Line \#5. The position
    of the 3rd Judge shall be located at a point in minimum 2 m from the
    corner of the Boundary Line \#8. In case of two Judges setting the
    position of the 1st Judge shall be located at a point in minimum 3 m
    from the outer line \#1 on the left of the mat and the 2nd Judge
    shall be located at a point in minimum 3m from of the outer line \#3
    on the right. The positions of Judges may be altered to facilitate
    media, broadcasting and/or sports presentation.
  \item
    Position of Recorder \& IVR: Position of Recorder \& IVR shall be
    located at a point 2 m from the Outer Line \#1. Position of Recorder
    may be altered to accommodate the environment of the venue and
    requirements from media broadcasting and/or sports presentation.
  \item
    Positions of Coaches: The position of the Coaches shall be marked at
    a point in minimum 2 m or more from the center point of the Outer
    Line of each contestant's side. Position of the coaches may be
    altered to accommodate the environment of the venue and requirements
    from media broadcasting and/or sports presentation.
  \item
    Position of Inspection desk: The position of the Inspection Desk
    shall be near the entrance of the Competition Area for the
    inspection of the contestants' protective equipment.
  \end{enumerate}
\item
  Technical and Environmental Requirements for Competitiion Venues

  \begin{itemize}
  \tightlist
  \item
    The minimum illuminance at the competition venue requred for
    broadcasted events shall be \textgreater=1,600 lux.
  \item
    The illuminance at the training venue shall range between 750 and
    900 lux.
  \item
    The temperature of the competition venue shall be maintained between
    17°C and 24°C.
  \item
    For G14 or higher-graded events, the minimum seating capacity of the
    competition venue shall be 4,000 seats.
  \end{itemize}
\end{enumerate}

\section{(Explanation \#1)}\label{explanation-1}

Elastic mat: The degree of elasticity and slipperiness of the mat must
be approved by the WT before the competition.

\section{(Explanation \#2)}\label{explanation-2}

Color: The color scheme of the mat's surface must avoid giving a harsh
reflection, or tiring the contestant's or spectator's eyesight. The
color scheme must also be appropriately matched to the contestant's
equipment, uniform and the surface of the Contest Area.

\section{(Explanation \#3)}\label{explanation-3}

Inspection Desk: At the Inspection desk, the inspector checks if all the
materials worn by the contestant are approved by the WT and fit the
contestant properly. In case they are found to be inappropriate, the
contestant is requested to change the protective equipment.

\chapter{Contestant}\label{contestant}

\begin{enumerate}
\item
  Qualification of Contestant

  \begin{enumerate}
  \item
    Holder of the nationality of the participating team
  \item
    One recommended by the WT MNA
  \item
    Holder of Taekwondo Dan/Poom certificate issued by the Kukkiwon
  \item
    Holder of the WT Global Athlete License (GAL)
  \item
    Contestants at the age of at least 17 years old for Senior in the
    year the pertinent tournament is held (15-17 years old for Junior
    Championships and 12-14 years old for Cadet Championships). Ages for
    Youth Olympic Games might be different depending on the decision of
    the IOC.
  \end{enumerate}
\end{enumerate}

\section{(Interpretation)}\label{interpretation-1}

The age limit is based on the year, not on the date. For instance, if
the Junior Championships are held on June 11, 2013, those contestants
born on between January 1, 1996 and December 31, 1998 are eligible to
participate.

\section{(Interpretation)}\label{interpretation-2}

Article 4.1 Qualification of Contestant: applied at WT promoted
championships, CU promoted championships, multi-sports games, and
approved event(s) of WT recognized international open taekwondo
tournament.

\section{(Interpretation)}\label{interpretation-3}

Article 4.1.1 \& 4.1.2 shall not be applied to WT recognized
international open taekwondo tournament.

\begin{enumerate}
\def\labelenumi{\arabic{enumi}.}
\setcounter{enumi}{1}
\item
  Contestant uniform and competition equipment

  \begin{enumerate}
  \item
    At the competitions enlisted in the WT Event Calendar, dobok or
    competition uniform and all competition equipment such as but not
    limited to mats, PSS, IVR and protective equipment must be those
    ones approved by the WT.

    \begin{enumerate}
    \tightlist
    \item
      Specifications of dobok or competition uniform, protective
      equipment, and all other equipment shall be set forth separately.
    \end{enumerate}
  \item
    A contestant shall wear a WT-approved dobok or competition uniform,
    trunk PSS, head PSS, groin guard, forearm guards, shin guards,
    gloves, sensing socks (in the case of using PSS) and be equipped
    with a mouthguard before entering the Field of Play. Head pss must
    be firmly tucked under left arms when entering into competition
    area. Head pss shall be put on the head following instructions of
    the referee before the start of the contest.
  \item
    In case of wearing dobok, the forearm and shin guards shall be worn
    beneath garment. In case of wearing competition uniform, the forearm
    and shin guards shall be worn in the garment or beneath garment. The
    groin guards shall be worn beneath garment for both cases.
  \item
    The contestant shall bring the WT-approved protective equipment, as
    well as gloves and the mouthguard, for his/her personal use. Wearing
    any item on the head other than the head pss shall not be permitted.
    The head protector for cadet athletes may be equipped with face
    shield. Any religious item shall be worn beneath the head pss and
    inside the dobok or competition uniform and shall not cause harm or
    obstruct the opposing contestant.
  \item
    Responsibilities of the organizing committee for competition
    equipment

    \begin{enumerate}
    \item
      The Organizing Committee of WT-promoted championships shall be
      responsible for preparing the following WT-recognized equipment
      for the use at the championships at its own expenses for all
      related materials, equipment and the related technicians for
      installation and operation.

      \begin{itemize}
      \tightlist
      \item
        Trunk PSS and Head PSS-related items and equipment - the choice
        of the PSS company shall be decided by the WT (For World
        Taekwondo Cadet Championships, conventional head protector with
        face shield shall be used)
      \item
        Mats
      \item
        Other protective equipment as reserve (Sensing socks, gloves,
        shin guards, forearm guards, groin guards and dobok or
        competition uniforms)
      \item
        Instant Video Replay (IVR) system and its related equipment,
        including but not limited to cameras (minimum 3 cameras per
        court and minimum 4 cameras, including one overhead camera for
        the semifinals and final); When broadcasting is available, the
        broadcast feed must be made available at the Video Replay desk
        for review purpose.
      \item
        Jumbo screen (for display of competition progress match tree,
        athlete profile. etc) inside the Field of Play (FOP)
      \item
        Spectator scoreboards (for display of instant video replay
        screen; minimum 12)
      \item
        Scoreboards at the court (for display of scoring; minimum 4 per
        court)
      \item
        Real Time Display System (RTDS) at athlete calling area and warm
        up area
      \item
        Real Time Referee Calling System (RTRCS) at referee lounge or
        waiting area.
      \item
        TV screen for showing competitions in live at referee lounge
      \item
        Metal detector at the inspection desk (minimum 2)
      \item
        Other competition equipment not prescribed in this article, if
        any, shall be described in competition's Operational Manual of
        the WT.
      \end{itemize}
    \item
      The Organizing Committee of WT-promoted championships shall be
      responsible for preparing the following equipment and materials,
      etc. at the training venue at its own expenses.

      \begin{itemize}
      \tightlist
      \item
        Trunk PSS and Head PSS-related items and equipment
      \item
        Mats
      \item
        Stationary Bicycle
      \item
        Running Machine
      \item
        Emergency equipment (refer to medical code for detailed
        information)
      \item
        Ice in the baskets
      \item
        Refrigerators
      \item
        Bottled water
      \end{itemize}
    \item
      It is the responsibility of the Organizing Committee to obtain
      approval of the WT on the number of the equipment to be prepared.
    \end{enumerate}
  \end{enumerate}
\item
  Anti-Doping Test

  \begin{enumerate}
  \item
    At the taekwondo events promoted or recognized by the WT, any use or
    administration of drugs or chemical substances described in the WADA
    Prohibited List is prohibited. The WADA Anti-doping Code shall be
    applied to the taekwondo competitions of the Olympic Games and other
    multi-sports Games. The WT Anti-Doping Rules shall be applied to WT
    promoted and/or recognized championships.
  \item
    The WT may carry out any doping tests deemed necessary to ascertain
    if a contestant has committed a breach of this rule, and any
    contestant who refuses to undergo this testing or who proves to have
    committed such a breach shall be removed from the final standings,
    and the record shall be transferred to the contestant next in line
    in the competition standings.
  \item
    The Organizing Committee shall be responsible for making all
    necessary preparations for conducting doping tests.
  \item
    The details of the WT Anti-Doping Rules shall be enacted as part of
    the bylaws.
  \end{enumerate}
\end{enumerate}

\section{(Explanation \#1)}\label{explanation-1-1}

Holder of the nationality of the participating team: When a contestant
is a representative of a national team, his/her nationality is decided
by citizenship of the country he/she is representing before submission
of the application to participate. Verification of citizenship is done
by inspection of the passport.\\
A competitor who is a national of two or more countries at the same time
may represent either one of them, as he/she may elect. However, in case
of changing nationality, he/she is allowed to represent other country
only if thirty-six (36) months have passed since competitor represented
a country in such events:

\begin{enumerate}
\def\labelenumi{\roman{enumi})}
\tightlist
\item
  Olympic Games
\item
  Qualification Tournaments for Olympic Games
\item
  4 year cycle continental multi sports games
\item
  2 year cycle continental championships
\item
  World Championships promoted by the WT
\end{enumerate}

This period may be reduced or even cancelled, with the agreement of the
NOCs and the WT. The WT may take disciplinary actions at any time
against the athlete and his/hers MNA that violates this article
including but not limited to deprival of the achievements.

\section{(Explanation \#2)}\label{explanation-2-1}

One recommended by the WT MNA: Each MNA is responsible for control of
non-pregnancy and gender and shall ensure that all team members have
been given medical exams that show them to be of adequate health and
fitness to participate. Also each MNA assumes full responsibilities for
accident and health insurance as well as the civil liabilities for their
contestants and officials during the WTpromoted championships.

\section{(Explanation \#3)}\label{explanation-3-1}

Mouthguard:\\
The color of the mouthguard is limited to white or transparent. It must
be at least 3mm thickness and cover entire upper teeth. Athletes with
dental braces needs to wear special mouthguard for braces, that covers
both upper and lower teeth, recommended by their dentist and submit the
letter from their dentist stating that the athlete is safe to compete
with the mouthguard that the dentist recommended.(Refer to WT
Mouthguard, Taping, Brace, Piercing rules for detailed information.)

\section{(Explanation \#4)}\label{explanation-4}

Head PSS: The color of head pss other than blue or red shall not be
permitted to compete.

\section{(Explanation \#5)}\label{explanation-5}

Instant Video Replay System: It is the responsibility of the Organizing
Committee to ensure broadcasting feed is provided for Instant Video
Replay for the matches requested by the WT.

\section{(Explanation \#6)}\label{explanation-6}

Taping: Taping of feet and hands will be strictly checked during the
athlete inspection process. The inspector may request the WT Commission
Doctor's approval for excessive taping. Contestants need to take off
taping on general weigh-in to see whether there is any open wound, cut
or bleeding. (Refer to WT Mouthguard, Taping, Brace, Piercing rules for
detailed information.)

\section{(Explanation \#7)}\label{explanation-7}

Any athlete who fails to wear appropriate safety protection gear and
equipment or remove potentially harmful material from their body at the
inspection as below (1) \textasciitilde{} (3) shall not be allow to
participate in the competition. (Refer to WT Mouthguard, Taping, Brace,
piercing rules and WT Medical Code Appendix III for detailed
information.)

\begin{enumerate}
\def\labelenumi{(\arabic{enumi})}
\tightlist
\item
  Any athlete whose protection gear (such as head, body, groin, hand,
  foot protector and mouthguard) has either inadequate body part
  coverage by the gear, inappropriate size, or significant defect of
  shape (or material).
\item
  Any athlete who does not have proper mouthguard per WT mouthguard
  rules at the inspection and does not have proper mouthguard at any
  time during the match, or fail to wear proper mouthguard despite of
  receiving a warning by an inspection referee, center referee or WT
  commissioned doctor.
\item
  Any athlete who has piercing, earing or any hard material in the face
  or any body part at the inspection or during the match despite of
  receiving a warning by an inspection referee, center referee or WT
  commissioned doctor.
\end{enumerate}

\section{(Explanation \#8)}\label{explanation-8}

Any athlete who participate in competition must have valid annual
periodic health evaluation (or medical certificate). Without valid
annual Medical Certificate shall not be allowed to participate in the
competition Medical Certificate. (Refer to WT medical code 8.3. Periodic
Health Evaluation and Appendix III A. Medical Certificate for detailed
information)

\section{(Explanation \#9)}\label{explanation-9}

Any athlete shall be disqualified for participation in competition if
the athlete does not follow the safety protection equipment rule, has
any health conditions that may jeopardized the safety by jurisdiction or
have disqualifying conditions in the medical certificate. (refer to WT
medical code Appendix III B. Disqualification for detailed information)

\section{(Explanation \#10)}\label{explanation-10}

Any athlete must have valid travel health insurance for the competition
to participate and submit the copy of the health insurance certificate
to GMS upon registration.

\chapter{Weight category}\label{weight-category}

ToDo: A táblázatok teljes elkészítése szükséges

\chapter{Classification and methods of
competition}\label{classification-and-methods-of-competition}

\begin{enumerate}
\item
  Competitions are classified as follows.

  \begin{enumerate}
  \item
    Individual competition shall normally be between contestants in the
    same weight category. When necessary, adjoining weight category may
    be combined to create a single classification. No contestant is
    allowed to participate in more than one (1) weight category in one
    event.
  \item
    Team Competition: Method and weight categories of team competition
    shall be stipulated in the Standing Procedures for World Taekwondo
    World Cup Team Championships.
  \end{enumerate}
\item
  Systems of competition are divided as follows.

  \begin{enumerate}
  \item
    Single elimination tournament system
  \item
    Round robin system
  \end{enumerate}
\item
  Taekwondo competition of the Olympic Games and 4 year cycle
  continental multi sports games may use single elimination tournament
  system or the combination of single elimination tournament system with
  repechage.
\item
  All international-level competitions recognized by the WT shall be
  formed with participation of at least four (4) countries in tournament
  with no less than four (4) contestants competed in each weight
  category, and any tournament with less than four (4) countries or
  weight category with less than 4 competed contestants cannot be
  recognized in the official results.
\item
  World Taekwondo Grand Prix Series will be organized based on the most
  recent Standing Procedure of World Taekwondo Grand Prix Series.
\end{enumerate}

\section{(Interpretation)}\label{interpretation-4}

\begin{enumerate}
\def\labelenumi{\arabic{enumi}.}
\item
  In the tournament system, competition is founded on an individual
  basis. However, the team standing can also be determined by the sum of
  the individual standings according to the overall scoring method.

  Team Standing

  Team standing shall be decided by the total points based on the
  following criteria.

  \begin{itemize}
  \tightlist
  \item
    Basic one (1) point per each contestant who entered the competition
    area after passing the general weigh-in
  \item
    One (1) point per each win (win by a bye included)
  \item
    Additional one hundred and twenty (120) points per gold medal
  \item
    Additional fifty (50) points per silver medal
  \item
    Additional twenty (20) points per bronze medal
  \end{itemize}

  In case more than two (2) teams are tied in score, the rank shall be
  decided by 1) number of gold, silver and bronze medals won by the team
  in order, 2) number of participating contestants and 3) higher points
  in heavier weight categories.
\item
  In the team competition system, the outcome of each team competition
  is determined by the individual team results.
\end{enumerate}

\section{(Explanation \#1)}\label{explanation-1-2}

Consolidated weight categories:

The method of consolidation shall follow the Olympic weight categories.

\chapter{Duration of Contest}\label{duration-of-contest}

\begin{enumerate}
\item
  Duration of Contest is classified as follows.

  \begin{enumerate}
  \item
    The duration of the contest shall be three rounds of two minutes
    each, with a one-minute rest period between rounds. In case of a tie
    score after the completion of the 3rd round, a 4th round of one
    minute will be conducted as the Golden round, after a one-minute
    rest period following the 3rd round.
  \item
    In the best of three (3) system, the duration of the contest shall
    be three rounds of two minutes each with a oneminute rest period
    between rounds. However, a 4 th round of one minute will not be
    conducted as the Golden round. In case of a tie score for
    corresponding round, the round winner shall be decided by the
    Article 15.
  \item
    In the World Cup Team Championships, the duration of the contest
    shall be three rounds of four (4) minutes of 1st round, five (5)
    minutes of 2nd and 3rd round with a one (1) minute rest period
    between rounds. The 1st round shall be conducted based on
    traditional team match format for one (1) minute per contest and the
    2nd and the 3rd round shall be conducted for five (5) minutes based
    on tag-team match format.
  \end{enumerate}
\item
  The duration of each round may be adjusted to 1 minute x 3 rounds, 1
  minute 30 seconds x 3 rounds, 2 minutes x 2 rounds or 5 minute x 1
  round (with 1 time out for 30 seconds to each contestant) upon the
  decision of the Technical Delegate for the pertinent championships.
\end{enumerate}

\chapter{Drawing of Lots}\label{drawing-of-lots}

\begin{enumerate}
\item
  The date of the drawing of lots shall be set forth in the outline of
  the championships. At least one representative from each team must
  attend the drawing of lots and participating teams are responsible for
  confirming their entries before the drawing of lots. In case of no
  representative can be at the drawing of lots, the team must designate
  a proxy and inform Technical delegate before the drawing of lots.
\item
  The drawing of lots may be conducted by random computerized drawing or
  by random manual drawing of lots. The method and order of drawing
  shall be determined by the Technical Delegate.
\item
  All ranked athletes will be seeded in all WT Promoted and Recognized
  tournaments unless otherwise stipulated in the relevant standing
  procedures or event outlines.
\end{enumerate}

\chapter{Weigh-in}\label{weigh-in}

\begin{enumerate}
\item
  The general weigh-in of the contestants on the day of competition
  shall be organized one day prior to the competition. The time for the
  general weigh-in will be decided by the Organizing Committee and be
  informed at the head of team meeting. The duration of the general
  weigh-in shall be two (2) hours at the maximum.
\item
  The random weigh-in will take place at the venue in the morning of the
  competition. All contestants who pass the general weigh-in must be
  present for random weigh-in maximum two (2) hours before the start of
  the competition. Should a contestant fail to appear for the random
  weigh-in, he/she will be disqualified. The random weigh-in must be
  completed maximum thirty (30) minutes before the start of the
  competition each day.

  \begin{enumerate}
  \tightlist
  \item
    The rate of selection for the random weigh-in shall be determined by
    the number of contestants in the weight category as following the
    criteria and subjects will be randomly selected by computerized
    system maximum two (2) hours before the start of the competition.
  \end{enumerate}

  \begin{enumerate}
  \def\labelenumii{\alph{enumii})}
  \tightlist
  \item
    More than 32 athletes : 20\% of total
  \item
    17-32 athletes : 6 athletes\\
  \item
    9-16 athletes : 4 athletes\\
  \item
    4-8 athletes : 2 athletes\\
  \item
    Below 4 athletes : None
  \end{enumerate}

  \begin{enumerate}
  \tightlist
  \item
    The random weigh-in shall be conducted with plus 5\% tolerance of
    the contestant's weight category. Underweight shall not be subjected
    to random weigh-in.
  \end{enumerate}
\item
  During the weigh-in, the male contestant shall wear underpants and the
  female contestant shall wear underpants and a brassiere. However,
  weigh-in may be conducted in the nude if the contestant wishes to do
  so.

  \begin{enumerate}
  \tightlist
  \item
    Cadet and junior contestant must be weighed with underwear(s) and
    100 grams will be allowed to compensate.
  \end{enumerate}
\item
  General weigh-in shall be made once, however, one more weigh-in is
  granted within the time limit to any contestant who did not qualify
  the first time. Random weigh-in shall be made only once per
  contestant, and there will not be 2nd weighin.
\item
  So as not to be disqualified during the weigh-in, scales identical to
  the official one shall be provided at the contestants' place of
  accommodation or at the competition venue for pre-weigh-in.
\end{enumerate}

\section{(Explantion\#1)}\label{explantion1}

The contestants on the day of competition:\\
This is defined as those contestants listed to compete on the scheduled
day by the Organizing Committee or the WT.

\section{(Explanation \#2)}\label{explanation-2-2}

A separate site for the weigh-in shall be installed for the male and
female contestants. The gender of weigh-in officials should be the same
as that of the contestants.

\section{(Explanation \#3)}\label{explanation-3-2}

Disqualification during the weigh-in:\\
When a contestant is disqualified at the weigh-in, the contestant shall
not be awarded any ranking points.

\section{(Explanation \#4)}\label{explanation-4-1}

Scales, identical to the official one:\\
The practice scale must be of the same type and calibrations as that of
the official scale and these facts must be verified prior to the
competition by the Organizing Committee.

Article 10 Procedure of the Contest

1

Call for contestants: The name of the contestants shall be announced at
the Athlete Calling Desk three (3) times beginning thirty (30) minutes
prior to the scheduled start of the contest. If a contestant fails to
report to the Desk following the third call, the contestant shall be
disqualified, and this disqualification shall be announced.

2

Inspection of body, uniform and apparatus: After being called, the
contestants shall undergo inspection of body, uniform and apparatus at
the designated inspection desk by the inspectors designated by the WT,
and the contestant shall not show any signs of aversion, and also shall
not wear any materials which may cause harm to the other contestant.

3

Entering the Competition Area: After inspection, the contestant shall
proceed to the Coach's zone with one coach and one team doctor or team
medical staff(such as physiotherapist, athletic trainer or chiropractor,
if any).

4

Procedure before the Beginning and After the End of the Contest

4.1 Before the start of the contest, the center referee will call
``Chung, Hong.'' Both contestants will enter the contest area with their
head pss firmly tucked under their left arms. When any of contestant is
not present or present without being fully attired, including all
protective equipment, uniform, etc., at the Coach's Zone by the time the
referee calls ``Chung, Hong'', he/she shall be regarded as withdrawn
from the contest and the referee shall declare the opponent as the
winner.

4.2 The contestants shall face each other and make a standing bow at the
referee's command of ``Cha-ryeot (attention)'' and ``Kyeong-rye (bow)''.
A standing bow shall be made from the natural standing posture of
``Charyeot'' by bending the waist at an angle of more than 30 degrees
with the head inclined to an angle of more than 45 degree. After the
bow, the contestants shall put on their head pss

4.3 The referee shall start the contest by commanding ``Joon-bi
(ready)'' and ``Shi-jak (start)''.

4.4 The contest in each round shall begin with the declaration of
``Shi-jak (start)'' by the referee.

4.5 The contest in each round shall end with the declaration of
``Keu-man (stop)'' by the referee. Even if the referee did not declare
``Keu-man'', the contest shall be deemed to have ended when the match
clock expired, however ``Gam-jeom'' can be given and registered in score
even after expiration of match clock.

4.6 The referee may pause a contest by declaring ``Kal-yeo'' (break) and
resume the contest by the command of ``Kye-sok'' (continue). When the
referee declares ``Kal-yeo'' the recorder should immediately stop the
match time; when the referee declares ``Kye-sok'' the recorder should
immediately restart the match time.

4.7 After the end of the last round, the referee shall declare the
winner by raising his/her hand to the winner's side.

4.7.1 In the best of three (3) system, the referee shall declare the
winner of respective round.

4.8

Retirement of the Contestants

5

Contest Procedure in Team Competition

5.1 Both teams shall stand facing each other in line in the submitted
team order towards the 1st Boundary Line from the Contestants' marks.

5.2 Procedure before the beginning and after the end of the contest
shall be conducted as in clause 4 of this article.

5.3 Both teams shall leave the Contest Area and stand by at the
designated area for each contestant's match.

5.4 Both teams shall line up in the Contest Area immediately after the
end of the final match facing each other.

5.5 The referee shall declare the winning team by raising his/her own
hand to the winning team's side.

(Explanation\#1) Team doctor, chiropractor, athlete trainer or a
physiotherapist: At the time of submission of entry for team officials,
copies of relevant and appropriate licenses of team doctor chiropractor,
athletic trainer or physiotherapist written in English shall be
attached. After verification, special accreditation cards shall be
issued to those team doctors, chiropractor, athlete trainer or
physiotherapists. Only those who have obtained proper accreditation
shall be allowed to proceed to competition area with coach. Only medical
doctor with active medical license who graduated from recognized medical
school can obtain accreditation card of team doctor. Other medical
staff(team chiropractor, athletic trainer, physiotherapist or other
healthcare professionals) are not allowed to claim themselves as team
doctor, which shall be regarded as improper accreditation.

(Guideline for officiating) In the case of using PSS, the referee shall
check if the PSS system and sensing socks worn by both athletes are
properly working. This process, however, may be deleted to save time for
speedy competition management.

Article 11

Permitted techniques and areas

1

Permitted techniques

1.1 Fist technique: A straight punching technique using the knuckle part
of a tightly clenched fist

1.2 Foot technique: Delivering techniques using any part of the foot
below the ankle bone

2

Permitted areas

2.1 Trunk: Attack by fist and foot techniques on the areas covered by
the trunk protector are permitted. However, such attacks shall not be
made on the part of the spine.

2.2 Head: The area above the collar bone. Only foot techniques are
permitted.

Article 12

Valid Points

1

Scoring Areas

1.1 Trunk: The blue or red colored area of the trunk protector

1.2 Head: The entire head above the bottom line of the head protector

2

Criteria for valid point(s):

2.1 Point(s) shall be awarded when a permitted technique is delivered to
the scoring areas of the trunk with a proper level of impact.

2.2 Point(s) shall be awarded when a permitted technique is delivered to
the scoring areas of the head.

2.3 The determination of the validity of the technique, level of impact,
and/or valid contact to the scoring area shall be made by the electronic
scoring system except fist techniques. These PSS determinations shall
not be subject to Instant Video Replay except for head kicks.

2.4 The WT Technical Committee shall determine the required level of
impact and sensitivity of the PSS, using different scales in
consideration of weight category, gender, and age groups. In certain
circumstances as deemed necessary the Technical Delegate may recalibrate
the valid level of impact.

3

The valid points are as follows.

3.1 One (1) point for a valid punch to the trunk protector

3.2 Two (2) points for a valid kick to the trunk protector

3.3 Four (4) points for a valid turning kick to the trunk protector

3.4 Three (3) points for a valid kick to the head

3.5 Five (5) points for a valid turning kick to the head

3.6 One (1) point awarded for every one ``Gam-jeom'' given to the
opponent contestant

4

Match score shall be the sum of points of the three rounds.

4.1 In the best of three (3) system, match score shall be the sum of the
number of round won of the three rounds.

5

Invalidation of point(s): When a contestant records points following
prohibited act(s):

5.1 If prohibited act is followed by point(s), the referee shall declare
the penalty for the prohibited act and invalidate the point(s).

(Explanation\#1)

Back kick (Dwichagi) is one type of turning kick techniques, the head
and shoulder rotation must occur to be considered as a back kick
(Dwichagi) and awarded technical points.

When contestants kick opponent by the back kick (Dwichagi), without
contestant's simultaneous rotation of head and shoulder, ``Back kick''
(Dwichagi) will not be considered turning kick.

Article 13

Scoring and publication

1

Scoring of valid point(s) shall be determined primarily using the
electronic scoring system installed in Protector and Scoring Systems
(PSS). Points awarded for punching techniques and additional points
awarded for turning kicks shall be scored by judges using manual scoring
devices. If PSS (Protector \& Scoring System) is not used, all scoring
shall be determined by judges using manual scoring devices.

2

If head PSS is not employed with trunk PSS, scoring for kicking
techniques to the head shall be made by judges using the manual scoring
devices.

3

The additional point given for a turning kick shall be invalidated if
the turning kick was not scored as a valid point(s) by PSS.

4

Under a three (3) corner judge setting, two or more judges shall be
needed to confirm valid scoring.

5

Under a two (2) corner judges setting, two judges shall be needed to
confirm valid scoring.

Article 14

Prohibited acts and Penalties

1

Penalties shall be declared by the referee.

2

Prohibited acts which described in article 14 shall be penalized with
``Gam-jeom'' by referee.

3

A ``Gam-jeom'' shall be counted as one (1) point for the opposing
contestant.

4

Prohibited acts 4.1 The following acts shall be classified as prohibited
acts, and ``Gam-jeom'' shall be declared. 4.1.1 Crossing the Boundary
Line 4.1.2 Falling down 4.1.3 Avoiding or delaying the match 4.1.4
Grabbing or pushing the opponent 4.1.5 The following are considered
prohibited acts: a) Lifting the leg to block b) Kicking the opponent's
leg to impede the opponent's kicking attack c) Kick was aiming to below
the waist d) Lifting the leg above waist for kicking in the air for four
(4) times or more e) Lifting a leg or kicking in the air for more than
three (3) seconds to impede opponent's potential attacking movements
4.1.6 Kicking below the waist 4.1.7 Attacking the opponent after
``Kal-yeo'' 4.1.8 Hitting the opponent's head with the hand 4.1.9
Butting or attacking with the knee 4.1.10 Attacking the fallen opponent
4.1.11 Attacking trunk PSS with the side or bottom of the foot in clinch
position 4.1.12 Attacking back of head PSS in clinch position 4.1.13
Following Misconducts of contestant or coach a) Not complying with the
referee's command or decision b) Inappropriate protesting behavior to
officials' decisions c) Inappropriate attempts to disturb or influence
the outcome of the match d) Provoking or insulting the opposing
contestant or coach e) Unaccredited doctor/physicians or other team
officials found to be seated in the doctor's position f) Any other
severe misconduct or unsportsmanlike conduct from a contestant or coach
g) When a contestant commits a prohibited Act followed by an Attack
After ``Kal-yeo'' (as per article 14.4.1.7) or any other unsportsmanlike
behavior (as per article 14.4.1.13) the Referee may give a 2nd
``Gam-jeom'' for `Attack after ``Kal-yeo''\,' or `Misconduct'.

4.2 When a coach or contestant commits excessive misconduct and does not
follow the referee's command the referee may declare a sanction request
by raising a yellow card. In this case the Competition Supervisory Board
shall investigate the contestant's and/or coach's behavior and determine
whether a sanction is appropriate

5

If a contestant intentionally and repeatedly refuses to comply with the
Competition Rules or the referee's orders, the referee may end the match
raising yellow card and declare the opposing contestant the winner.

6

If the referee at the inspection desk or officials in the Field of Play
determines, in consultation with the PSS technician, if necessary, that
a contestant or coach has attempted to manipulate the sensitivity of PSS
sensor(s) and/or inappropriately alter the PSS so as to affect its
performance, the contestant shall be disqualified.

7

When a contestant receives ten (10) ``Gam-jeom'', the referee shall
declare the contestant loser by referee's punitive declaration (PUN)

7.1 In the best of three (3) system, when a contestant receives five (5)
'' Gam-jeom'' in a round, the opponent will be declared the winner of
that round.

8

In Article 14.7, ``Gam-jeom'' shall be counted in the total score of the
three rounds.

(Interpretation) Objectives in establishing the prohibited acts and
penalties are as follows, (1) To secure the contestant's safety (2) To
ensure fair competition (3) To encourage appropriate techniques

(Explanation \#1) ``Gam-jeom''

\begin{enumerate}
\def\labelenumi{\roman{enumi}.}
\item
  Crossing the Boundary Line: A ``Gam-jeom'' shall be declared when one
  foot of a contestant crosses the Boundary Line. No ``Gam-jeom'' will
  be declared if a contestant crosses the boundary line as a result of a
  prohibited act by the opposing contestant.
\item
  Falling down: ``Gam-jeom'' shall be declared for falling down.
  However, if a contestant falls down due to the opponent's prohibited
  acts ``Gam-jeom'' penalty shall not be given to the fallen contestant,
  while a penalty shall be given to the opponent. If both contestants
  fall as a result of incidental collision, or in case a contestant who
  received a point with turning kick falls down, no penalty shall be
  given.
\item
  Avoiding or delaying the match:
\end{enumerate}

\begin{enumerate}
\def\labelenumi{\alph{enumi})}
\item
  This act involves stalling with no intention of attacking. A
  contestant who continuously displays a non-engaging style shall be
  given a ``Gam-jeom''. If both contestants remain inactive after three
  (3) seconds, the center referee will signal the ``\,``Gong-gyeok''
  command. A ``Gam-jeom'' will be declared: On both contestants if there
  is no activity from them three (3) seconds after the command was
  given; or on the contestant who moved backwards from the original
  position three (3) seconds after the command was given.
\item
  Turning the back and move away to avoid the opponent's attack should
  be punished as it expresses the lack of a spirit of fair play and may
  cause serious injury. The same penalty should also be given for
  evading the opponent's attack by bending below waist level or
  crouching.
\item
  Retreating from the technical engagement only to avoid the opponent's
  attack and to run out the clock, ``Gam-jeom'' shall be given to the
  passive contestant.
\item
  Pretending injury means exaggerating injury or indicating pain in a
  body part not subjected to a blow for the purpose of demonstrating the
  opponent's actions as a violation, and also exaggerating pain for the
  purpose of elapsing the match time.
\item
  ``Gam-jeom'' shall also be given to the athlete who asks the referee
  to stop the contest in order to adjust the position/fit of protective
  equipment.
\end{enumerate}

\begin{enumerate}
\def\labelenumi{\roman{enumi}.}
\setcounter{enumi}{3}
\tightlist
\item
  Grabbing or pushing the opponent:
\end{enumerate}

\begin{enumerate}
\def\labelenumi{\alph{enumi})}
\tightlist
\item
  This includes grabbing any part of the opponent's body, uniform or
  protective equipment with the hands. It also Includes the act of
  grabbing the foot or let or hooking the leg with forearm. For pushing,
  it is permitted as a quick impact and a contestant must disengage from
  opponent after one push. The flowing acts shall be penalized.
\end{enumerate}

\begin{itemize}
\tightlist
\item
  Pushing the opponent with prolonged or continuous contact
\item
  Pushing the opponent out of the boundary line
\item
  Pushing the opponent in a way that prevents kicking motion or any
  normal execution of attacking movement
\end{itemize}

\begin{enumerate}
\def\labelenumi{\alph{enumi}.}
\setcounter{enumi}{21}
\tightlist
\item
  Lifting the leg or cut kick motion shall not be penalized only when it
  is followed by execution of punching or kicking technique in
  combination motion.
\end{enumerate}

\begin{enumerate}
\def\labelenumi{\roman{enumi}.}
\setcounter{enumi}{5}
\item
  Attacking below the waist: This action applies to an attack on any
  part below the waist. When an attack below the waist is caused by the
  recipient in the course of an exchange of techniques, no penalty will
  be given. This article also applies to strong kicking or stamping
  actions to any part of the thigh, knee or shin for the purpose of
  interfering with the opponent's technique.
\item
  Attacking the opponent after ``Kal-yeo'':
\end{enumerate}

a)Attacking after Kal-yeo requires that the attack results in actual
contact to the opponent's body. b)If the attacking motion started before
the Kal-yeo, the attack shall not be penalized. c)In Instant Video
Replay, the timing of Kal-yeo shall be defined as the moment that the
referee's Kal-yeo hand signal was completed (with fully extended arm);
and the start the attack shall be defined as the moment that the
attacking foot is fully off the floor. d)If an attack after Kal-yeo did
not land on the opponent's body but appeared deliberate and malicious
the referee may penalize the behavior with a ``Gam-jeom''

\begin{enumerate}
\def\labelenumi{\roman{enumi}.}
\setcounter{enumi}{7}
\item
  Hitting the opponent's head with the hand: This article includes
  hitting the opponent's head with the hand (fist), wrist, arm, or
  elbow. However, unavoidable actions due to the opponent's carelessness
  such as excessively lowering the head or carelessly turning the body
  cannot be punished by this article.
\item
  Butting or attacking with the knee: This article relates to an
  intentional butting or attacking with the knee when in close proximity
  to the opponent. However, contact with the knee that happens in the
  following situations cannot be punished by this article. When the
  opponent rushes in abruptly at the moment a kick is being executed
  Inadvertently, or as the result of a discrepancy in distance in
  attacking.
\item
  Attacking the fallen opponent: This action is extremely dangerous due
  to the high probability of injury to the opponent. The danger arises
  from the following:

  \begin{itemize}
  \tightlist
  \item
    The fallen opponent is in an immediate defenseless state
  \item
    The impact of any technique which strikes a fallen contestant will
    be greater due to the contestant's position. These types of
    aggressive actions toward a fallen opponent are not in accordance
    with the spirit of taekwondo and as such are not appropriate to
    taekwondo competition. In this regard, penalties should be given for
    intentionally attacking the fallen opponent regardless of the degree
    of impact
  \end{itemize}
\end{enumerate}

When misconduct is committed by a contestant or a coach during a rest
period, past the five (5) seconds of the round conclusion, the referee
can immediately declare the ``Gam-jeom'' and the ``Gam-jeom'' shall be
recorded to the upcoming round. However, ``Gam-jeom'' shall be recorded
to the previous round if the action happened within five (5) seconds of
the round conclusion.

Article 15

Golden Points and Decision of Superiority

1

In the event the winner cannot be decided after 3 rounds, a 4th round
(golden round) will be conducted in one-minute round.

2

In case of a contest advances to a golden round, all scores awarded
during the first three (3) rounds shall be void.

3

The first contestant scores two (2) or more points or whose opponent
receives two ``Gam-jeoms'' in the golden round shall be declared the
winner.

4

In the event that neither contestant has scored two (2) points after the
completion of the golden round, the winner shall be decided by
superiority based the following criteria:

4.1 The contestant who received a point by a punch in the golden round
4.2 If none of the contestant received a point by a punch or both
contestants received a point by a punch each in the golden round, the
contestant who achieved a higher number of hits registered by the PSS
during the golden round. 4.3 If number of hits registered by the PSS is
tied, the contestant who won more rounds in first three rounds 4.4 If
number of round won is tied, the contestant who received less numbers of
a ``Gam-jeom'' during all four rounds 4.5 If the three above criteria
are the same, the referee and judges shall determine superiority based
on the content of the golden round. If the superiority decision is tied
among the referee and judges, the referee shall decide the winner.

5

In the best of three (3) system, in case of tie score for corresponding
round, the round winner shall be decided by superiority based on
following criteria:

5.1 Most points scored by turning or spinning kick. 5.2 If the technical
score is the same, the contestant who has more scored in the order of a
higher value techniques as follows (Head, Trunk, Punch, Gam-Jeom). 5.3
If the high value points are same, the contestant who received higher
number of hits registered by the PSS.. 5.4 If the three above criteria
are the same, the referee and judges shall determine superiority. a) In
case of two (2) corner judges, winner shall be decided by referee and
two (2) judges b) In case of three (3) corner judges, winner shall be
decided by three (3) judges except referee.

(Explanation \#1) Decision of superiority by judges shall be based on
technical dominance of an opponent through aggressive match management,
the greater number of techniques executed, the use of the more advanced
techniques both in difficulty and complexity, and display of the better
competition manner. (Explanation \#2) In the event that referee is
counting for delivering successful head kick which caused knock down
right before the opponent's body kick but only the latest body kick was
registered, the coach of the athlete who delivered head kick may request
for video replay. If review jury determines that head kick was valid and
performed earlier than body kick, the referee shall invalidate the
point(s) scored by body kick, then declare points for head kick, and
declare the one who delivered head kick as winner. In the event that one
athlete who already has one (1) point is delivering successful punch
right before the opponent's body kick but only the latest body kick was
registered, the coach of the athlete who delivered punch may request for
video replay. If review jury determines that punch was performed earlier
than body kick, the referee shall ask whether corner judges were scored
or not. If two or more corner judges scored for punch, referee shall
invalidate the point(s) scored by body kick, then declare points for
punch, and declare the one who delivered punch as winner

(Guideline for officiating) The procedure for superiority decision shall
be as follows except for the best of three (3) system. 1) Prior to the
contest, all refereeing officials take the superiority card with them.
2) When a match is to be decided by superiority, the referee shall
declare ``Woo-se-girok (Record Superiority)''. 3) Upon the referee's
declaration, the judges shall record the winner within 10 seconds with
their heads down, sign on the card and then give it to the referee. 4)
The referee shall collect all superiority cards, record the final
result, and then declare the winner. 5) Upon declaration of winner, the
referee shall hand over the cards to the recorder and the recorder shall
submit the cards to the Technical Delegate of World Taekwondo.

(Guideline for officiating for the best of three (3) system) 1) When a
round is to be decided by superiority, the referee shall declare
``Woo-se-girok (Record Superiority)''. 2) Upon the referee's
declaration, the judges declare the winner simultaneously after the
Referee's count of three (3) by using round winner hand signal facing
the head table. 3) In case of two (2) corner judges, winner shall be
decided by referee and two (2) corner judges. 4) In case of three (3)
corner judges, winner shall be decided by three (3) corner judges except
referee. 5) The Review Jury shall record the final result, and declare
to Operator the round winner. 6) Upon declaration of winner, the referee
shall declare the round or match winner.

Article 16

Decisions

1

Win by Referee Stops Contest (RSC)

2

Win by final score (PTF)

3

Win by point gap (PTG)

4

Win by golden points (GDP)

5

Win by superiority (SUP)

6

Win by withdrawal (WDR)

7

Win by disqualification (DSQ)

8

Win by referee's punitive declaration (PUN)

9

Win by disqualification for unsportsmanlike behavior (DQB)

(Explanation \#1)

Referee Stops Contest: The referee declares RSC in the following
situations:

\begin{enumerate}
\def\labelenumi{\roman{enumi}.}
\tightlist
\item
  if a contestant has been knocked down by an opponent's legitimate
  technique and cannot resume the contest by the count of ``Yeo-dul'';
  or if the referee determines the contestant is not able to resume the
  competition regardless of the progress of counting;
\item
  if a contestant disregards the referee's command to continue the match
  three times; iii.if the referee recognizes the need to stop the match
  to protect a contestant's safety;
\item
  when the Commission Doctor determines that the match should be stopped
  due to a contestant's injury;
\end{enumerate}

(Explanation \#2)

Win by point gap: In case of twenty (20) points difference between two
athletes at the time of the completion of 2nd round and/or at any time
during the 3rd round, the referee shall stop the contest and shall
declare the winner by point gap. Win by point gap shall not be applied
in semi- finals \& finals in senior division by the outline of
tournament.

(Explanation \#3)

Win by withdrawal: The winner is determined by withdrawal of the
opponent. - When a contestant withdraws from the match due to injury or
other reasons - When the coach throws a towel into the court to signify
forfeiture of the match

(Explanation \#4)

Win by disqualification: This is the result determined by the
contestant's failure in weigh-in or when a contestant fails to report to
the Athlete Calling Desk following the third call.

The follow-up actions should be different in accordance with the reason
of disqualification.

\begin{enumerate}
\def\labelenumi{\roman{enumi}.}
\item
  In the event that contestants have not passed or did not show at
  weigh-in: The result should be reflected on the draw sheet and the
  information should be provided to technical officials and all relevant
  persons. Referees will not be assigned for this match. The opponent of
  athletes that did not pass or show at weigh-in will not need to appear
  at the court to compete.
\item
  In the event that a contestant passed weigh-in but did not appear at
  the Athletes Calling Desk: The assigned referee and opponent shall
  enter the FOP and waiting in their position until the referee declares
  the opponent a winner of the match. Detailed procedure is stipulated
  in 4.1 of Article 10.
\end{enumerate}

(Explanation \#5)

Win by the referee's punitive declarations:

The referee declares PUN in the following situations:

\begin{enumerate}
\def\labelenumi{\roman{enumi}.}
\tightlist
\item
  if a contestant accumulated ten (10) ``Gam-jeom''
\end{enumerate}

(Explanation \#6) Win by disqualification for unsportsmanlike behavior:
DQB shall be declared in the following situations: - When a contestant
or any of his/her team member is found manipulating of the sensor(s) or
scoring system of the PSS - When a contestant cheats the process of
weigh-in - When a contestant is found violating the WT Anti-Doping Rules
- When a contestant or coach commits serious infringing behavior
described in article 23.3.1\& 23.3.2

The all result of contestant who lost by DQB shall be removed, and other
contestant's result which affected by DQB shall be reallocated.

(Explanation \#7)

In the best of three (3) system, the decisions shall follow the
procedure of Article 16: - 16.1 Win by Referee Stops Contest (RSC) -
16.2 Win by Final Score (PTF) - 16.6 Win by Withdrawal (WDR) - 16.7 Win
by Disqualification (DSQ) - 16.9 Win by Disqualification for
unsportsmanlike behavior (DQB)

\begin{enumerate}
\def\labelenumi{\roman{enumi})}
\tightlist
\item
  In case of Article 16.2. Win by final score (PTF), match score shall
  be the sum of the number of round won of the three rounds.
\end{enumerate}

ii)In case of Article 16.3. Win by point (PTG), in case of twelve (12)
points difference between two athletes per round, the referee shall stop
the contest and shall declare the winner by point gap for corresponding
round. Point gap for corresponding round shall not be applied in semi-
finals \& finals in senior division by the outline of tournament.

Article 17

Knock Down

A Knock Down shall be declared, when a legitimate attack is delivered
and;

1

When any part of the body other than the sole of the foot touches the
floor due to the force of the opponent's scoring technique

2

When a contestant is staggered and shows no intention or ability to
continue as a result of the opponent's scoring techniques.

3

When the referee judges that the contest cannot continue as the result
of being struck by a legitimate scoring technique

(Explanation \#1)

A knock down:

This is the situation in which a contestant is knocked to the floor or
is staggered or unable to respond adequately to the requirements of the
match due to a blow. Even in the absence of these indications, the
referee may interpret as a knock down, the situation where, as the
result of contact, it would be dangerous to continue or when there is
any question about the safety of a contestant.

Article 18

Procedure in the event of a Knock Down

1

When a contestant is knocked down as the result of the opponent's
legitimate attack, the referee shall take the following measures.

1.1 The referee shall keep the attacker away from downed contestant by
declaration of ``Kal-yeo (break)''. The recorder shall stop the match
clock following the referee's ``Kal-yeo (break)'' command. 1.2 The
referee shall first check the status of the downed contestant and count
aloud from ``Ha-nah (one)'' up to ``Yeol (ten)'' at one second intervals
towards the downed contestant, making hand signals indicating the
passage of time. 1.3 In case the downed contestant stands up during the
referee's count and desires to continue the fight, the referee shall
continue the count up to ``Yeo-dul (eight)'' for recovery of the
contestant. The referee shall then determine if the contestant is
recovered and, if so, continue the contest by declaration of ``Kye-sok
(continue)''. 1.4 When a contestant who has been knocked down cannot
demonstrate the will to resume the contest by the count of ``Yeo-dul
(eight)'', the referee shall announce the other contestant winner by
RSC(Referee Stops Contest). 1.5 In case both contestants are knocked
down, the referee shall continue counting as long as one of the
contestants has not sufficiently recovered. 1.6 In case both contestants
are knocked down and both contestants fail to recover by the count of
``Yeol'', the winner shall be decided by the match score before the
occurrence of Knock Down. 1.7 When it is judged by the referee that a
contestant is unable to continue, the referee may decide the winner
either without counting or during the counting.

2

Procedures to be followed after the contest: Any contestant who could
not continue the match as a result of a serious injury regardless of any
parts of the body cannot enter competition within thirty (30) days
without approval of the WT Medical Chairman after submission of a
statement from the physician designated by the pertinent national
federation.

2.1 Except for medical emergency, any contestant with any serious injury
must be evaluated by venue medical doctor and confirmed by medical
chairman (MC) at medical room immediately after the contest. 2.2 Any
contestant who had knockout due to head injury must be checked by
medical doctor at medical room per WT medical rules. A venue medical
doctor must perform SCAT5 on the injured contestant for diagnosis of
concussion in case of head injury within 30minutes after the head
injury. 2.3 Any significant(moderate to severe) head trauma or
concussion carries mandatory suspension for any competition during the
suspension period(see 18.2.5) This mandatory medical suspension period
cannot be shortened in any circumstances once the suspension is given.
2.4 The decision on the suspension of the contestant for significant
head trauma or concussion must be made based the on one of the following
criteria (1)\textasciitilde(3). (1) Comprehensive neurological
examination and neurocognitive testing (SCAT 5 or other validated
concussionassessment tools permitted by MC Chair) performed by
commissioned doctor in the venue medical room which must be reported to
MC Chair) (2) Referee-stop-contest due to any loss of consciousness,
altered mental status or inability to make a meaningful, stable and
voluntary movement as a result of a direct head trauma at least for ten
(10) seconds or by the count ten (10) (3) Failure to fully recover from
head trauma and resume the match within one (1) minute of medical
evaluation on the mat after centre referee calls a doctor for possible
concussion. 2.5 Any contestant who had diagnosis of significant head
trauma or concussion based on one of the above 2.4 criteria will get
30days suspension(senior), 40days suspension(junior) or 50days
suspension(cadet). This mandatory medical suspension period cannot be
shortened in any circumstances once the suspension is given. 2.6 Any
contestant who had second concussion in last 90 days will get 90days
suspension and who had third concussion in last 180days will get 180days
suspension. 2.7 For any incidence of concussion or significant head
trauma underreported, misdiagnosed or mismanaged without mandatory
medical suspension, WT Sport Department and WT MC Medical Committee led
by WT MC Chair shall start investigating the incidence by retrospective
video review even after the end of the competition. The incidence must
be reported to WT MC Chair within 30days after the incidence date to
commence the investigation. If the video review confirmed by at least
three (3) reviewers of WT medical committee reveals obvious concussion
or serious head trauma (knockout more than 10 seconds) or other serious
injuries which mandate at least 30-day mandatory medical suspension, WT
medical committee shall override the medical examiner (OMD or
commissioned doctor)'s decision and apply the mandatory suspension rules
to the athlete to protect the health and safety of the contestant.

(Explanation \#1)

Keep the attacker away: In this situation the standing opponent shall
return to the respective contestant's mark, however, if the downed
contestant is on or near the opponent's contestant's mark, the opponent
shall wait at the boundary line in front of his/her coach's chair.

(Guideline for officiating)

The referee must be constantly prepared for the sudden occurrence of a
knock down or situation where the contestant is staggered, which is
usually characterized by a powerful blow accompanied by impact.

(Explanation \#2)

In case the downed contestant stands up during the referee's count and
desires to continue the fight: The primary purpose of counting is to
protect the contestant. Even if the contestant desires to continue the
match before the count of eight is reached, the referee must count until
``Yeo-dul (eight)'' before resuming the match. Counting to ``Yeo-dul''
is compulsory and cannot be altered by the referee.

*Count from one to ten: Ha-nah, Duhl, Seht, Neht, Da-seot, Yeo-seot,
Il-gop, Yeo-dul, A-hop, Yeol.

(Explanation \#3)

The referee shall then determine if the contestant has recovered and, if
so, restart the contest by the declaration of ``Kye-sok'': The referee
must ascertain the ability of the contestant to continue while he/she
counts until eight. Final confirmation of the contestant's condition
after the count of eight is only procedural and the referee must not
needlessly pass time before resuming the contest.

(Explanation \#4)

When a contestant who has been knocked down cannot express the will to
resume by the count of ``Yeo-dul'', the referee shall announce the other
contestant winner by RSC after counting to ``Yeol'': The contestant
expresses the will to continue the match by gesturing several times in a
fighting position with the clenched fists. If the contestant cannot
display this gesture by the count of ``Yeo-dul'', the referee must
declare the other contestant winner after first counting ``A-hop'' and
``Yeol''. Expressing the will to continue after the count of ``Yeo-dul''
cannot be considered valid. Even if the contestant expresses the will to
resume by the count of ``Yeo-dul'', the referee can continue counting
and may declare the contest over if he/she determines the contestant is
incapable of resuming the match.

(Explanation \#5)

When a contestant is downed by a powerful scoring blow and whose
condition appears serious, the referee can suspend the count and call
for first aid or do so in conjunction with the count.

(Guideline for officiating)

\begin{enumerate}
\def\labelenumi{\roman{enumi}.}
\item
  The referee must not spend additional time confirming the contestant's
  recovery after counting to ``Yeo-dul'' as a result of failing to
  observe that condition during the administration of the count.
\item
  When the contestant clearly recovers before the count of ``Yeo-dul''
  and expresses the will to resume and the referee can clearly discern
  the contestant's condition yet resumption is hampered by the
  requirement of medical treatment, the referee must first resume the
  match with the declaration of ``Kye-sok'' and immediately after
  declare ``Kal-yeo'' and ``Kye-shi'' and then follow the procedures of
  Article 19.
\end{enumerate}

Article 19

Procedures of suspending the match

1

When a contest is to be stopped due to the injury of one or both
contestants, the referee shall take the measures prescribed below.
However, in a situation which warrants suspending the contest for
reasons other than an injury, the referee shall declare''Kal-yeo
(break)'' and resume the contest by declaring ``Kye-sok (continue)''.

1.1 The referee shall suspend the contest by declaration of ``Kal-yeo''
and order the recorders to suspend the time.

1.2 The referee shall allow the contestant one minute to receive first
aid by the commission doctor; the referee may allow team doctor to treat
first aid if the commission doctor is not available or if it is deemed
necessary.

1.2.1 The commission doctor may request more time (up to 2 minutes) if
necessary.

1.2.2 If there is no commission doctor, team doctor or medical chairman
available, any doctor (or medical associate) near competition mat can be
requested to provide the athlete with first aid.

1.3 If an injured contestant cannot return to the match after one minute
the referee shall declare the other contestant winner.

1.4 In case resumption of the contest is impossible after one minute,
the contestant causing the injury by a prohibited act to be penalized by
``Gam-jeom'' shall be declared the loser.

1.5 In case both the contestants are knocked down and are unable to
continue the contest after one minute, the winner shall be decided upon
points scored before the injuries occurred.

1.5.1 In case of best of three (3) System: If both contestants are
knocked down and are unable to continue the contest after one minute at
the Round 1 or Round 3, the winner shall be determined by the points
scored before the injuries occurred in the pertinent Round. If this
occurs during the Round 2, the winner shall be determined by the
decision of the Round 1.

1.5.2 If points are tied, the winner shall be decided according to
criteria of superiority.

1.6 If the referee determines a contestant's pain is caused only by a
bruise the referee shall declare ``Kal-yeo'' and give a command to
resume the match with the call, ``stand-up''. If the contestant refuses
to continue the match after the referee gives the command ``stand up''
three times, the referee shall declare the match `Referee Stops
Contest'.

1.7 If the referee determines a contestant has received an injury such
as broken bone(s), dislocation, sprain ankle(s), and/or bleeding, the
referee shall allow the contestant to receive a first aid treatment for
one minute after ``Kyeshi''. The referee may allow the contestant to
receive first aid treatment even after giving the commanding ``standup''
if the contestant is determined to be injured in one of the categories
above.

1.8 Stopping the match due to injury: If the referee determines a
contestant has received an injury such as broken bone(s), dislocation,
sprain ankle(s), and/or bleeding, the referee shall consult with the
chairperson of the Medical Committee or the commissioned doctor assigned
by the chairperson. If a contestant is re-injured in the same manner,
the chairperson of the Medical Committee or the commission doctor
assigned by the chairman may advise the referee to stop the match and
declare the injured the loser.

(Explanation \#1)

When the referee determines that the competition cannot be continued due
to injury or any other emergency situation, he/she may take the
following measures:

\begin{enumerate}
\def\labelenumi{\roman{enumi}.}
\item
  If the situation is critical such as a contestant losing consciousness
  or suffering from a severe injury and time is crucial, first aid must
  be immediately directed first and the match must be closed. In this
  case, the result of the match will be decided as follows.

  \begin{itemize}
  \tightlist
  \item
    The contestant causing the injury shall be declared the loser if the
    outcome was the result of a prohibited act to be penalized by
    ``Gam-jeom''.
  \item
    The incapacitated contestant shall be declared the loser if the
    outcome was the result of a legal action or accidental, unavoidable
    contact.
  \item
    If the outcome was unrelated to the match contents, the winner shall
    be decided by the match score before suspension of the match. If the
    suspension occurs before the end of the first round, the match shall
    be invalidated. If points are tied, the winner shall be decided
    according to the criteria of superiority.
  \item
    In case of best of three (3) System: If the outcome was unrelated to
    the match contents, the winner shall be determined by the points
    scored before the suspension occurred in the pertinent Round (in
    case of the Round 1 or 3). If points are tied, the winner shall be
    decided according to the criteria of superiority. If this occurs in
    the Round 2, the winner shall be determined by the decision of the
    Round 1.
  \end{itemize}
\item
  If first aid treatment is need for an injury, the contestant can
  receive necessary treatment within one minute after the declaration of
  ``Kye-shi''.
\end{enumerate}

\begin{verbatim}
- Order to resume the match: It is the decision of the center referee, after consultation with the Commission Doctor, whether or not it is possible for the contestant to resume the match. The referee can anytime order the contestant to resume the match within one minute. The referee can declare any contestant who does not follow the order to resume the match the loser of the contest.
- While the contestant is receiving medical treatment or is in the process of recovering, 40 seconds after the declaration of "Kye-shi", the referee begins to loudly announce the passage of time in five second intervals.  When the contestant cannot return to the Contestant's Mark by the end of the one minute period, the match results must be declared.
- After the declaration of "Kye-shi", the one minute time interval must be counted from the moment the commission doctor enter the mat or after waiting for commission doctor up to 10 seconds if not readily available in the mat. However, when the doctor's treatment is required but the doctor is absent or additional treatment is necessary, the one minute time limit can be suspended by the judgment of the referee.
- If resumption of the match is impossible after one minute, the decision of the match will be determined according to sub-article "i" of this article.
\end{verbatim}

\begin{enumerate}
\def\labelenumi{\roman{enumi}.}
\setcounter{enumi}{2}
\tightlist
\item
  If both contestants become incapacitated and are unable to resume the
  match after one minute or urgent conditions arise, the match result is
  decided according to the following criteria:
\end{enumerate}

\begin{verbatim}
- If the outcome is the result of a prohibited act to be penalized by "Gam-jeom" by one contestant that person shall be the loser.
- If the outcome was not related to any prohibited act to be penalized by "Gam-jeom", the result of the match shall be determined by the match score at the time of suspension of the match. However, if the suspension occurs before the end of the first round, the match shall be invalidated and the Organizing Committee will determine an appropriate time to re-contest the match. The contestant who cannot resume the match shall be deemed to have withdrawn from the match.
- In case of best of three (3) System: If the outcome was unrelated to "Gam-jeom" penalty, the contest result shall be determined by the points scored before the suspension occurred in the pertinent Round (in case of Round 1 or 3). If the points are tied, the winner shall be determined according to the criteria of superiority. If this occurs in the Round 2, the winner shall be determined by the decision of the Round 1.
- If the outcome is the result of prohibited acts to be penalized by "Gam-jeom" by both contestants, then both contestants shall lose.
\end{verbatim}

(Explanation \#2)

The situation which warrants suspending the match beyond the
above-prescribed procedures shall be treated as follows.

\begin{enumerate}
\def\labelenumi{\roman{enumi}.}
\item
  When uncontrollable circumstances require suspension of the match, the
  referee shall suspend the match and follow the directives of the
  Technical Delegate.
\item
  If the match is suspended after the completion of the second round,
  the outcome shall be determined according to the match score at the
  time of suspension.
\end{enumerate}

\begin{verbatim}
- In case of best of three (3) System, if the match is suspended, the outcome shall be determined by the point scored before the suspension occurred in pertinent Round in case of Round 1 or Round 3 of the contest. If points are tied, the winner shall be decided according to the criteria of superiority. In case of the occurrence during Round 2, the winner shall be decided upon decision of Round 1.
\end{verbatim}

\begin{enumerate}
\def\labelenumi{\roman{enumi}.}
\setcounter{enumi}{2}
\tightlist
\item
  If the match is suspended before the conclusion of the second round, a
  rematch shall, in principle, be conducted and shall be held in three
  rounds.
\end{enumerate}

Article 20

Technical Officials

1

Technical Delegate (TD)

1.1 Qualification: WT President shall appoint TD among members of WT
technical committee for WT promoted championships upon recommendation of
WT Secretary General.

1.2 Roles: TD is responsible to ensure that WT Competition Rules are
properly applied and preside over the Head of Team meeting and drawing
of lots session. TD approves the result of draw, weigh-in and
competitions before it being officialized. TD has the right to make
final decisions on competition area and overall technical matters on
competitions in consultation with Competition Supervisory Board. TD
shall make final decisions on any matters pertaining to competitions not
prescribed in Competition Rules. TD serves as the Chairman of
Competition Supervisory Board. TD is responsible for reporting of event
evaluation.

2

Competition Supervisory Board (CSB) Member

2.1 Qualification: CSB members shall be appointed by the WT President
upon recommendation of Secretary General from those who have sufficient
experience and knowledge of taekwondo competitions.

2.2 Composition: CSB shall consist of one Chairperson and no more than 4
members at WT-promoted championships. Chairpersons of WT Games
Committee, WT Referee Committee, and WT Medical Committee and WT Athlete
Committee shall be included in CSB as ex-officio members. The
composition, however, may be adjusted by the President, if necessary.

2.3 Roles: CSB shall assist TD in competitions and technical matters and
ensure the competitions are held in accordance with the schedule. CSB
shall evaluate the performances of Review Jury and refereeing officials.
CSB shall also concurrently act as the Extraordinary Sanctions Committee
during competition with regard to competition management matters.

3

Refereeing officials

3.1 Qualification: Holders of International Referee Certificate
registered by the WT

3.2 Duties

3.2.1 Referee

3.2.1.1 The referee shall have control over the match.

3.2.1.2 The referee shall declare ``Shi-jak'', ``Keu-man'', ``Kal-yeo'',
``Kye-sok'', ``Kye-shi'', ``Shi-gan'', winner and loser, deduction of
points, penalty, and retiring. All the referees' declarations shall be
made after the results are confirmed.

3.2.1.3 The referee shall have the right to make decisions independently
in accordance with the prescribed rules.

3.2.1.4 In principle, the center referee shall not award points.
However, if one of the corner judges raises his/her hand because a point
was not scored, then the center referee will convene a meeting with the
judges. If it was found that two corner judges request for change of the
judgment, the referee must accept and correct the judgment (in case of 1
referee + 3 judges). In a two corner judge setting, the result of the
scoring can be revised when two persons among two judges and the referee
agree to do so.

3.2.1.5 In case as defined by the Article 15, the decision of
superiority shall be made by refereeing officials after the end of four
(4) rounds when necessary.

3.2.2 Judges

3.2.2.1 The judges shall mark the valid points immediately.

3.2.2.2 The judges shall state their opinions forthrightly when
requested to do so by the referee.

3.2.3 Review Jury (RJ)

3.2.3.1 RJ shall review an instant replay and inform the referee of the
decision within thirty (30) seconds.

3.2.4 Technical Assistant

3.2.4.1 TA shall keep monitoring scoreboard during the contest if the
scoring, penalties and timing are correctly publicized, and immediately
notify the referee of any problematic issue in this regard.

3.2.4.2 TA shall notify the referee of starting or stopping the contest
in close communication with system operator and recorder.

3.2.4.3 TA manually records all scores, penalties and IVR result in TA
paper.

3.3 Composition of refereeing officials per court

3.3.1 The officials' squad is composed of one (1) referee and three (3)
judges.

3.3.2 The officials' squad is composed of one (1) referee and two (2)
judges

3.4 Assignment of refereeing officials

3.4.1 The assignment of the referees and judges shall be made after the
contest schedule is fixed. 3.4.2 Referees and judges with the same
nationality as that of either contestant shall not be assigned to such a
contest. However, an exception can be made for the judges when the
number of refereeing officials is insufficient.

3.5 Responsibilities for decisions: Decisions made by the referees and
judges shall be conclusive and they shall be responsible to the
Competition Supervisory Board for the content of those decisions.

3.6 Uniforms

3.6.1 The referees and judges shall wear the uniform designated by the
WT.

3.6.2 The refereeing officials shall not carry or take any materials to
the contest area which might interfere with the contest. Use of mobile
phones by refereeing officials in the field of play may be restricted,
if necessary.

4 Recorders: The recorder shall time the contest, periods of time-out,
and suspensions, and also shall record and publicize the awarded points,
and/or penalties.

(Explanation \#1)

Refereeing officials must stay in a separate hotel to avoid any contact
with team officials. The hotel should be located less than 20 minutes
distance by car from the venue.

(Interpretation)

The details of the refereeing official's qualifications, duties,
organization, etc. shall follow the WT Regulations on the Administration
of International Referees.

(Interpretation) TD may replace or penalize the refereeing officials in
consultation with CSB in the event that refereeing officials have been
wrongly assigned, or when it is judged that any of the assigned
refereeing officials have unfairly conducted the contest or made
unjustifiable mistakes repeatedly.

(Guideline for officiating)

In case that each judge awards different score respectively to the legal
attack on the head, for instance, one judge gives one point, another
gives two and the other gives no point, and that no point is recognized
as a valid one, or in the case that the recorder makes mistakes in
timing, scoring or penalties, any of the judges may indicate the mistake
and ask for confirmation among the judges. Then, the referee may declare
``Kal-yeo'' (break)'' to stop the contest and gather the judges to ask
for statements. After discussion, the referee must publicize the
resolution. In the case that a coach requests for video review for the
same case that one of the judges requests for a meeting between
refereeing officials, the referee shall first gather judges before
taking the request from coach. If it has been decided to correct the
decision, the coach shall remain seated without using appeal quota. If
the coach still stands and request for video review, the referee shall
take the coach's request.

Article 21

Instant Video Replay

\begin{enumerate}
\def\labelenumi{\arabic{enumi}.}
\tightlist
\item
\end{enumerate}

In case there is an objection to a judgment of the refereeing officials
during the contest, the coach of a team can make a request to the center
referee for an immediate review of the video replay. The coach can only
request video replay for followings;

\begin{enumerate}
\def\labelenumi{\roman{enumi})}
\tightlist
\item
  Penalties against the opponent for instances of falling down or
  crossing the boundary line or attacking the opponent after ``Kal-yeo''
  or attacking the fallen opponent
\item
  Technical point
\item
  Any penalty against own contestant
\item
  Any mechanical malfunction or error in time management. In case of
  appeal for PSS mechanical malfunction, the coach may request to the
  center referee for a testing of the PSS at any time during the 2 nd
  and/or the 3rd round. However, if the PSS mechanical function is
  properly working, the coach's appeal quota shall be forfeited. As well
  as this coach's appeal is considered to be a misconduct of coach that
  ``Gam-jeom'' shall be given to own contestant in in accordance with
  Article 14.4.1.3 `Following Misconducts of contestant of coach'. This
  only applies to the `Best of 3' System.
\item
  When referee forgot to invalidate point(s) after ``Gam-jeom'' was
  given for prohibited act
\item
  Wrong identification of fist attacking contestant by judge
\item
  Head kick that is not scored
\end{enumerate}

\begin{enumerate}
\def\labelenumi{\arabic{enumi}.}
\setcounter{enumi}{1}
\tightlist
\item
\end{enumerate}

When coach appeals, the center referee will approach the coach and ask
the reason for the appeal. Any appeal shall not be admissible on any
points scored by foot or fist attacks on the trunk or foot attack on
trunk PSS. Coach may request instant video replay for head kick
regardless of using head PSS. The scope of instant video replay request
is limited to the only one action which has occurred within five (5)
seconds from the moment of the coach's request. Once the coach rises the
blue or red card to request for instant video replay, it will be
considered that the coach has used his/her allocated appeal under any
circumstance unless the judge's meeting satisfies the coach.

3

Referee shall request the Review Jury to review the instant video
replay. Review Jury, who is not of the same nationality as the
contestants, shall review the video replay.

3.1 In the last five (5) seconds of any Round, the Center Referee may
request for IVR review to check the possible Gam-jeom penalties for the
following actions:

\begin{verbatim}
- Falling down
- Crossing the boundary line
- Attack after Kal-yeo
- Attack the fallen opponent
\end{verbatim}

\begin{itemize}
\tightlist
\item
  Any points scored after Prohibited Act will be invalidated.
\end{itemize}

3.2 If a referee perceives a contestant to be staggering, a strong
impact to the head, kick to the eye(s), bleeding or knocked down by a
kick to the head, and so begins counting, but the attack was not scored
by the head PSS, the referee or coach must request IVR to make the
decision for awarding or not awarding points after the count.

3.3 Referee may request IVR for clarification before declaration of
``Gam-jeom'' for pretending injury.

4 After review of the instant video replay, the Review Jury shall inform
the center referee of the final decision within thirty (30) seconds
after receiving the request.

5 Coach shall be allocated with one (1) appeal to request an instant
video replay request per each contest. However, based on the size and
level of the Championships, the Technical Delegate may decide the number
of appeal quota during the head of team meeting. If the appeal is
successful and the contested request is corrected, the coach shall
retain the appeal right for the pertinent contest.

6 The decision of the Review Jury is final; no further appeals during
the contest or protest after the contest will be accepted except errors
in the calculation of scoring in contest of Olympic Qualification
Tournaments, G6 event or above as explained article 21.7.1.

7 In the case that there is a clear erroneous decision from the
refereeing officials and scoring operators on identification of the
contestant or errors in the calculation of scoring as follows

\begin{verbatim}
- Score and Gam-jeom input by Operator
- Misidentification of athletes by Center Referee
\end{verbatim}

any of the refereeing officials shall request for IVR and correct the
decision at any time during the round.

7.1 If errors in calculation of scoring or misidentification of the
athlete were not corrected during the round by the refereeing officials
and the errors subsequently affected the result of the winner of the
contest or the round in Olympic Qualification Tournaments, G6 events or
above, the Technical Delegate (TD) has the right to review the case with
assistance of the Competition Supervisory Board (CSB) to make the
necessary correction within 30 minutes after the pertinent round. If the
errors were identified and confirmed, the TD may request for resumption
of the round from the moment when the error was occurred. In case of
multiple errors, competition will be resumed from the moment when the
first error was occurred.

8 In the case of a successful appeal, the Competition Supervisory Board
may investigate the contest at the end of the competition day and take
disciplinary action against the concerned refereeing officials, if
necessary.

9 In any time during the round, any of the judges can ask for adding or
removing technical points regardless of coach's appeal quota.

10 In the tournament where instant video replay system is not available,
the following protest procedure will be applied.

10.1 In case there is an objection to a referee judgment, an official
delegate of the team must submit an application for re-evaluation of
decision (protest application) together with the non-refundable protest
fee of US\$200 to the Board of Arbitration (Competition Supervisory
Board) within 10 minutes after pertinent contest.

10.2 Deliberation of re-evaluation shall be carried out excluding those
members with the same nationality as that of contestant concerned, and
resolution on deliberation shall be made by majority.

10.3 The members of the Board of Arbitration (Competition Supervisory
Board) may summon the refereeing officials for confirmation of events.

10.4 The resolution made by the Board of Arbitration (Competition
Supervisory Board) will be final and no further means of appeal will be
applied.

10.5 Deliberation procedures are as follows:

10.5.1 A coach or head of team from the protesting nation shall be
permitted to make a brief verbal presentation to the Board of
Arbitration in support of their position. The coach or head of team from
the respondent nation shall be allowed to present a brief rebuttal.

10.5.2 After reviewing the protest application, the contest of the
protest must be arranged according to the criterion of ``Acceptable'' or
``Unacceptable''. 10.5.3 If necessary, the Board can hear opinions from
the referee or judges.

10.5.4 If necessary, the Board can review the material evidence of the
decision, such as the written or visual recorded data.

10.5.5 After deliberation, the Board shall hold the secret ballot to
determine a majority decision.

10.5.6 The Chairperson will make a report documenting the outcome of the
deliberation and shall make this outcome publicly known.

10.5.7 Subsequent process following the decision:

10.5.7.1 Errors in determining the match results, mistakes in
calculating the match score or misidentifying a contestant shall result
in the decision being reversed.

10.5.7.2 Error in application of the rules: When it is determined by the
Board that the referee made a clear error in applying the Competition
Rules, the outcome of the error shall be corrected and the referee shall
be punished.

10.5.7.3 Errors in factual judgment: When the Board decides that there
was a clear error in judging the facts such as impact of striking,
severity of action or conduct, intention, timing of an act in relation
to a declaration or area, the decision shall not be changed and the
officials seen to have made the error shall be reprimanded.

Article 22

Deaf-Taekwondo

This article outlines the modifications to the Competition Rules used
for Deaf-Taekwondo. For matters not covered by Article 22 the WT
Competition Rules will apply.

1 Qualification of athlete Contestant must have gone through
classification procedures as outlined in the World Para-Taekwondo and
DeafTaekwondo Classification Code and been assigned Sport Class and
Sport Class Status

2 Weight Categories Olympic weight categories apply to competitions in
Deaf-Taekwondo

3 World Deaf-Taekwondo Championships will be organized based on the most
recent Standing Procedure of the World Deaf-Taekwondo Championships.

Article 23

Sanctions

1 The WT President, Secretary General, or Technical Delegate may request
that an on-spot Extraordinary Sanctions Committee be convened for
deliberation when inappropriate behaviors may have been committed by a
coach, a contestant, an official, and/or any member of a Member National
Association.

2 The Extraordinary Sanctions Committee shall investigate the matter,
and summon person(s) concerned for confirmation of events.

3 The Extraordinary Sanctions Committee shall deliberate the matter and
determine if disciplinary actions are to be imposed. The result of
deliberation shall be immediately announced to the public. If there is a
finding of violation, a written decision, including the relevant facts,
rules, supporting evidence (such as witness statements), the sanction
imposed, and rationale, shall be given to the sanctioned party as soon
as reasonably practicable, and a copy shall be included in the Technical
Delegate's report.

3.1 Potential violations on Conduct of a Contestant;

3.1.1 Refusing the referee's command to complete the ending procedures
of the match, including not participating in the declaration of the
winner.

3.1.2 Throwing his/her belongings (head protector, gloves, etc.) as an
expression of dissatisfaction with decision.

3.1.3 Not leaving the competition area after the end of a match

3.1.4 Not returning to a match after a referee's repeated command

3.1.5 Not complying with the competition official's ruling or command

3.1.6 Not complying with the Competition Management Officials reasonable
instructions related to the orderly management of the event

3.1.7 Manipulation of scoring equipment, sensors or/and any part of a
PSS

3.1.8 Any serious unsportsmanlike behavior during a match or aggressive
misconduct toward competition officials.

3.2 Potential violations on conduct of a coach, team official, or other
members of a Member National Association;

3.2.1 Complaining about or/and arguing against an official's decision
during or after a round.

3.2.2 Arguing with the referee or other official(s)

3.2.3 Violent behavior or remark toward officials, opponents or the
opposing side, or spectators during a match

3.2.4 Provoking spectators or spreading false rumor

3.2.5 Instructing athlete(s) to participate in misconduct, such as
remaining in the competition area after a match

3.2.6 Violent behaviors such as throwing or kicking personal
belonging(s) or competition material(s).

3.2.7 Not following instructions of competition officials to leave the
Field of Play or Venue

3.2.8 Any other serious misconducts toward competition officials

3.2.9 Any attempt to bribe competition officials

4 Disciplinary actions: Disciplinary actions issued by the Extraordinary
Sanction Committee may vary according to the degree of the violation.
The following sanctions may be given:

4.1 Disqualification of the athlete

4.2 Warning and order to issue official apology

4.3 Removal of accreditation

4.4 Ban from the competition venue

\begin{enumerate}
\def\labelenumi{\roman{enumi})}
\item
  Ban for the day
\item
  Ban for the duration of the Championships.
\end{enumerate}

4.5 Cancellation of result

\begin{enumerate}
\def\labelenumi{\roman{enumi})}
\item
  Cancellation of the match result and all related merits
\item
  Cancellation of WT Ranking points
\end{enumerate}

4.6 Monetary fine of between \$100-to-\$5,000 US dollars per violation.

5 The Extraordinary Sanctions Committee may recommend to the WT or WT at
its own initiative may investigate and determine, that additional
disciplinary actions be taken against the members/MNA involved,
including but not limited to longer-term suspension, lifetime ban,
and/or additional monetary fines. Such recommendation can be based on
violations of the Competition Rules and Interpretations as well as
violations of the WT Code of Ethics or other pertinent WT rules.

Article 24

Other matters not specified in Competition Rules

1 In the case that any matters not specified in the rules occur, they
shall be dealt with as follows.

1.1 Matters related to a contest shall be decided through consensus by
the refereeing officials of the pertinent contest.

1.2 Matters not related to a specific contest throughout the
championships such as technical matters, competition matters, etc. shall
be decided by the Technical Delegate.

\end{document}
